\chapter{Formula Collection}

This section provides a compact overview of all central formulae of this work.

% ============================================================
\section{Quadratic Sequence}

\textbf{Main function:}
\begin{equation}
k(n) = \frac{1}{6}(2n^2 + 3n + 1)
\end{equation}

\textbf{Congruence condition:}
\begin{equation}
k(n) \in \mathbb{Z}
\quad\Longleftrightarrow\quad
n \equiv 1,5 \pmod{6}
\end{equation}

\textbf{$y$-intercept:}
\begin{equation}
k(0) = \frac{1}{6}
\end{equation}

% ============================================================
\section{Difference Sequence}

\textbf{Generator function} (Theorem~5.4, Ch.~\ref{sec:spirale}):
\begin{equation}
a_k = \frac{1}{2}\bigl[(12k+3) + (-1)^{k+1}(4k+1)\bigr], \qquad k \geq 1.
\end{equation}
The sequence begins with $a_1 = 10$ (major), $a_2 = 9$ (minor), $a_3 = 26$, \ldots

\textbf{Major steps} (odd $k = 2m-1$):
\begin{equation}
a_{2m-1} = \Delta_{\mathrm{Major}}(m) = 16m - 6
\end{equation}

\textbf{Minor steps} (even $k = 2m$):
\begin{equation}
a_{2m} = \Delta_{\mathrm{Minor}}(m) = 8m + 1
\end{equation}

% ============================================================
\section{Circles and Radii}

\textbf{Circle circumferences:}
\begin{equation}
U_m = a_{2m-1} + a_{2m} = 24m - 5
\end{equation}

\textbf{Constant difference:}
\begin{equation}
\Delta U = 24
\end{equation}

\textbf{Radius sequence:}
\begin{equation}
R_m = \frac{U_m}{2\pi} = \frac{24m - 5}{2\pi}
\end{equation}

\textbf{Radial difference (key result):}
\begin{equation}
\boxed{
\Delta R = \frac{12}{\pi}
= 3.819718634205\ldots
}
\end{equation}

% ============================================================
\section{Archimedean Spiral}

\textbf{Standard form:}
\begin{equation}
r(\theta) = a\theta,
\qquad
a = \frac{6}{\pi^2}
\end{equation}

\textbf{Spiral equation:}
\begin{equation}
\boxed{
r(\theta) = \frac{6}{\pi^2}\theta
}
\end{equation}

\textbf{Connection to the Basel problem:}
\begin{equation}
a = \frac{6}{\pi^2}
= \frac{1}{\zeta(2)},
\qquad
\zeta(2) = \sum_{n=1}^{\infty}\frac{1}{n^2} = \frac{\pi^2}{6}
\end{equation}

% ============================================================
\section{Simplified Parametrisation}

\textbf{Polar form with parameter $p$:}
\begin{equation}
\boxed{
\begin{aligned}
r(p) &= \frac{6}{\pi\sqrt{3}}\sqrt{p}, \\[4pt]
\theta(p) &= \frac{\pi}{\sqrt{3}}\sqrt{p}
\end{aligned}
}
\end{equation}

\textbf{Cartesian coordinates:}
\begin{equation}
\begin{aligned}
x(p) &= \frac{6}{\pi\sqrt{3}}\sqrt{p}
       \cos\!\left(\frac{\pi}{\sqrt{3}}\sqrt{p}\right), \\[4pt]
y(p) &= \frac{6}{\pi\sqrt{3}}\sqrt{p}
       \sin\!\left(\frac{\pi}{\sqrt{3}}\sqrt{p}\right)
\end{aligned}
\end{equation}

\textbf{Derivatives:}
\begin{equation}
\begin{aligned}
\frac{dr}{dp} &= \frac{\sqrt{3}}{\pi\sqrt{p}}, \\[4pt]
\frac{d\theta}{dp} &= \frac{\pi}{2\sqrt{3}\sqrt{p}}
\end{aligned}
\end{equation}

% ============================================================
\section{Arc-Length Integral}

\textbf{Exact form:}
\begin{equation}
\boxed{
L = \int_{p_1}^{p_2}
\frac{1}{\pi}\sqrt{\frac{3}{p} + \pi^2}\, dp
}
\end{equation}

\textbf{Alternative form:}
\begin{equation}
L = \int_{p_1}^{p_2}
\sqrt{1 + \frac{3}{\pi^2 p}}\, dp
\end{equation}

\textbf{Taylor approximation:}
\begin{equation}
\boxed{
L \approx (p_2 - p_1)
+ \frac{3}{2\pi^2}\ln\!\left(\frac{p_2}{p_1}\right)
+ \frac{9}{8\pi^4}\left(\frac{1}{p_1} - \frac{1}{p_2}\right)
}
\end{equation}

\textbf{Coefficients:}
\begin{equation}
\begin{aligned}
\frac{3}{2\pi^2} &= 0.151981775464\ldots, \\[4pt]
\frac{9}{8\pi^4} &= 0.011549230037\ldots
\end{aligned}
\end{equation}

% ============================================================
\section{Conical Helix}

\textbf{Cone equation:}
\begin{equation}
z = -\frac{3}{4} + \frac{\pi}{2}r
\end{equation}

\textbf{Height function:}
\begin{equation}
z(\theta) = -\frac{3}{4} + \frac{3}{\pi}\theta
\end{equation}

\textbf{Half-angle:}
\begin{equation}
\tan(\alpha) = \frac{2}{\pi}
\quad\Rightarrow\quad
\boxed{
\alpha = \arctan\!\left(\frac{2}{\pi}\right)
= 32.481636590530^\circ
}
\end{equation}

\textbf{Helix parametrisation:}
\begin{equation}
\boxed{
\gamma(\theta) =
\begin{pmatrix}
\dfrac{6}{\pi^2}\theta \cos\theta \\[6pt]
\dfrac{6}{\pi^2}\theta \sin\theta \\[6pt]
-\dfrac{3}{4} + \dfrac{3}{\pi}\theta
\end{pmatrix}
}
\end{equation}

\textbf{Integer points:}
\begin{equation}
\theta_n = \frac{\pi}{3}n + \frac{\pi}{4}
\qquad (n \equiv 1,5 \pmod{6})
\end{equation}

\textbf{Tangent vector:}
\begin{equation}
\gamma'(\theta) =
\begin{pmatrix}
\dfrac{6}{\pi^2}(\cos\theta - \theta\sin\theta) \\[6pt]
\dfrac{6}{\pi^2}(\sin\theta + \theta\cos\theta) \\[6pt]
\dfrac{3}{\pi}
\end{pmatrix}
\end{equation}

\textbf{Planar curvature (approximation):}
\begin{equation}
\kappa_{\mathrm{2D}}(\theta)
= \frac{\theta^2 + 2}{a\,(\theta^2 + 1)^{3/2}},
\qquad
a = \frac{6}{\pi^2}
\end{equation}

\textbf{Exact 3D curvature and torsion} (Theorem~\ref{thm:3d-curvature-torsion}):
\begin{align}
\kappa_{\mathrm{3D}}(\theta)
  &= \frac{a\,\sqrt{\,c_1^2(\theta^2+4) + a^2(\theta^2+2)^2\,}}
          {\bigl(a^2(1+\theta^2) + c_1^2\bigr)^{3/2}},\\[4pt]
\tau(\theta)
  &= \frac{c_1\,(\theta^2 + 6)}
          {c_1^2\,(\theta^2+4) + a^2\,(\theta^2+2)^2},
\qquad c_1 = \frac{3}{\pi}.
\end{align}

\clearpage

% ============================================================
\section{Numerical Constants}

\begin{table}[h]
\centering
\caption{Important constants (15 decimal places)}
\begin{tabular}{lcl}
\toprule
\textbf{Constant} & \textbf{Value} & \textbf{Meaning} \\
\midrule
$\pi$ & 3.141592653589793 & Circle constant \\
$\pi^2$ & 9.869604401089358 & \\
$\sqrt{3}$ & 1.732050807568877 & \\
$\dfrac{6}{\pi^2}$ & 0.607927101854027 & Spiral parameter \\
$\dfrac{12}{\pi}$ & 3.819718634205488 & Radial difference \\
$\dfrac{6}{\pi\sqrt{3}}$ & 1.102657790843584 & Radius factor \\
$\dfrac{\pi}{\sqrt{3}}$ & 1.813799364234218 & Angular factor \\
$\dfrac{3}{2\pi^2}$ & 0.151981775464066 & Taylor coefficient 1 \\
$\dfrac{9}{8\pi^4}$ & 0.011549230036520 & Taylor coefficient 2 \\
$\arctan(\frac{2}{\pi})$ & $32.481636590530^\circ$ & Cone half-angle \\
\bottomrule
\end{tabular}
\end{table}
